\subsection{POINCARÉ SERIES OF GRADED ALGEBRAS}

Let K be a commutative ring with unit element, not reduced to 0. Let M be a graded K-module of type Z, and \(M_{n}\) the set of homogeneous elements of M of degree n.~Assume that each \(M_{n}\) is free and of finite type. Then the rank \(\operatorname{rk}_{\mathbf{K}}(\mathbf{M}_{n})\) is defined for all n (Commutative Algebra, Chap. II, §5, no. 3).

DEFINITION 1. If there exists \(n_0 \in \mathbf{Z}\) such that \(M_n = 0\) for \(n \leqslant n_0\) , the formal series \(\sum_{n \geqslant n_0} \mathrm{rk}_K(M_n) T^n\) , which is an element of \(\mathbf{Q}((T))\) , is called the Poincaré series of \(M\) and denoted by \(P_M(T)\) .

Let \(M'\) be another graded K-module of type Z, and \((M_n')_{n \in \mathbf{Z}}\) its grading. Assume that \(M_n'\) is zero for n less than a certain number. Then

\[
\mathrm{P} _ {\mathrm{M} \oplus \mathrm{M} ^ {\prime}} (\mathrm{T}) = \mathrm{P} _ {\mathrm{M}} (\mathrm{T}) + \mathrm{P} _ {\mathrm{M} ^ {\prime}} (\mathrm{T})\tag{1}
\]

and, if \(\mathbf{M} \otimes_{\mathbf{K}} \mathbf{M}'\) is provided with the total grading (Algebra, Chap. II, §11, no. 5),

\[
\mathrm{P} _ {\mathrm{M} \otimes \mathrm{M} ^ {\prime}} (\mathrm{T}) = \mathrm{P} _ {\mathrm{M}} (\mathrm{T}) \mathrm{P} _ {\mathrm{M} ^ {\prime}} (\mathrm{T}).\tag{2}
\]

PROPOSITION 1. Let \(S = \bigoplus_{n \geqslant 0} S_n\) be a commutative graded K-algebra with a system of generators \((x_1, x_2, \ldots, x_m)\) consisting of homogeneous and algebraically independent elements. Let \(d_i\) be the degree of \(x_i\) , and assume that \(d_i > 0\) for all \(i\) . Then the \(S_n\) are free and of finite rank over \(K\) , and

\[
\mathrm{P} _ {\mathrm{S}} (\mathrm{T}) = \prod_ {i = 1} ^ {m} (1 - \mathrm{T} ^ {d _ {i}}) ^ {- 1}.\tag{3}
\]

Indeed, S can be identified with the tensor product \(K[x_{1}] \otimes \cdots \otimes K[x_{m}]\) , provided with the total grading. The Poincaré series of \(K[x_{i}]\) is

\[
\sum_ {n \geqslant 0} \mathrm{T} ^ {n d _ {i}} = (1 - \mathrm{T} ^ {d _ {i}}) ^ {- 1},
\]

and it suffices to apply (2).

Under the assumptions of Prop. 1, we shall say that S is a graded polynomial algebra over K.

COROLLARY. The degrees \(d_{i}\) are determined up to order by S.

Indeed, the inverse of \(\mathrm{P}_{\mathrm{S}}(\mathrm{T})\) is the polynomial \(\mathrm{N}(\mathrm{T}) = \prod_{i=1}(1 - \mathrm{T}^{d_i})\) , which is thus uniquely determined. If q is an integer \(\geqslant 1\) and if \(\zeta \in C\) is a primitive qth root of unity, the multiplicity of the root \(\zeta\) of \(\mathrm{N}(\mathrm{T})\) is equal to the number of \(d_{i}\) that are multiples of q. This number is zero for q sufficiently large. The number of \(d_{i}\) equal to q is thus determined uniquely by descending induction.

The integers \(d_{i}\) are called the characteristic degrees of S. The number of them is equal to the transcendence degree of S over K when K is a field; we shall also call it the transcendence degree of S over K in the general case. It is the multiplicity of the root 1 of the polynomial N(T).

Let \(S = \bigoplus_{n \geqslant 0} S_n\) be a commutative graded K-algebra, and \(R = \bigoplus_{n \geqslant 0} R_n\) a graded subalgebra of S. Assume that each \(R_n\) is free and of finite type, and that the R-module S admits a finite basis consisting of homogeneous elements \(z_1, z_2, \ldots, z_N\) of degrees \(f_1, f_2, \ldots, f_N\) . Then, if M denotes the graded K-module \(\sum_{j=1}^{N} Kz_j\) , the graded K-module S is isomorphic to \(R \otimes_K M\) , so each \(S_n\) is free and of finite type and

\[
P _ {S} (T) = P _ {M} (T) P _ {R} (T) = \left(\sum_ {j = 1} ^ {N} T ^ {f _ {j}}\right) P _ {R} (T).\tag{4}
\]

PROPOSITION 2. Retain the preceding notation and assume that \(S\) and \(R\) are graded polynomial algebras.

\begin{enumerate}
\def\labelenumi{(\roman{enumi})}
\item
  R and S have the same transcendence degree \(r\) over \(K\) .
\item
  Let \(p_1, \ldots, p_r\) (resp. \(q_1, \ldots, q_r\) ) be the characteristic degrees of \(S\) (resp. \(R\) ). Then
\end{enumerate}

\[
\prod_ {i = 1} ^ {r} \left(1 - \mathrm{T} ^ {q _ {i}}\right) = \left(\sum_ {j = 1} ^ {\mathrm{N}} \mathrm{T} ^ {f _ {j}}\right) \prod_ {i = 1} ^ {r} \left(1 - \mathrm{T} ^ {p _ {i}}\right).
\]

\begin{enumerate}
\def\labelenumi{(\roman{enumi})}
\setcounter{enumi}{2}
\tightlist
\item
  \(Np_{1}p_{2}\ldots p_{r}=q_{1}q_{2}\ldots q_{r}.\)
\end{enumerate}

Formula (4) shows first of all that the multiplicity of the root 1 is the same in the polynomials \(\mathrm{P}_{\mathrm{S}}(\mathrm{T})^{-1}\) and \(\mathrm{P}_{\mathrm{R}}(\mathrm{T})^{-1}\) and, taking (3) into account, proves both (i) and (ii).

It follows from (ii) that

\[
\prod_ {i = 1} ^ {r} \left(1 + \mathrm{T} + \mathrm{T} ^ {2} + \dots + \mathrm{T} ^ {q _ {i} - 1}\right) = \left(\sum_ {j = 1} ^ {\mathrm{N}} \mathrm{T} ^ {f _ {j}}\right) \prod_ {i = 1} ^ {r} \left(1 + \mathrm{T} + \mathrm{T} ^ {2} + \dots + \mathrm{T} ^ {p _ {i} - 1}\right).
\]

Putting \(\mathrm{T} = 1\) in this relation gives (iii).

Remark. Let \(S = K[X_1, \ldots, X_n]\) be a graded polynomial algebra over \(K\) , \(d_i\) the degree of \(X_i\) and \(F(X_1, \ldots, X_n)\) a homogeneous element of degree \(m\) of \(S\) . Then

\[
\sum_ {i = 1} ^ {n} d _ {i} X _ {i} \frac {\partial F}{\partial X _ {i}} = m F.\tag{5}
\]

Indeed, it is immediate that the K-linear map D from S to S that transforms every homogeneous element z of degree p into pz is a derivation of S. Thus

\[
m \mathrm{F} (\mathrm{X} _ {1}, \dots , \mathrm{X} _ {n}) = \mathrm{D} (\mathrm{F} (\mathrm{X} _ {1}, \dots , \mathrm{X} _ {n})) = \sum_ {i = 1} ^ {n} \mathrm{D} (\mathrm{X} _ {i}) \frac {\partial \mathrm{F}}{\partial \mathrm{X} _ {i}} = \sum_ {i = 1} ^ {n} d _ {i} \mathrm{X} _ {i} \frac {\partial \mathrm{F}}{\partial \mathrm{X} _ {i}}.
\]

